% ===================================================================
%  TABLICA Nr 7 — Wieś zimą. — Zagroda wiejska wiosną.
%  Meme regime que les precedents. Premiere de la DEUXIEME SERIE :
%  l'apparat porte « SERIA DRUGA », ordinal apres le nom, comme au
%  tableau 1 — c'est le membre polonais de SERIO qui le lit.
%
%  LE MALE, LA FEMELLE ET LE PETIT SONT TROIS MOTS, ET L'IDO LES
%  FABRIQUE : kavalo / kavalino / kavalyuno, bovulo / bovino /
%  bovyuno, mutonulo / mutonino / mutonyuno. Le polonais a des mots
%  pleins — koń, klacz, źrebię ; byk, krowa, cielę ; baran, owca,
%  jagnię ; kozioł, koza, koźlę — et chacun porte son renvoi. Une
%  colonne qui dirait « młody koń » pour źrebię perdrait le mot que la
%  planche montre.
%
%  ET « pastoro » EST LE BERGER, NON LE PASTEUR. Le faux ami est
%  d'autant plus facile que la scene precedente parle d'un cure ; ici
%  c'est « pasterz », avec son chien et sa houlette.
% ===================================================================

%%K t07-noto-1 noto
\VUnotes{143.2mm}{%
(*) Zobacz co do szczegółów posiłku tablice 6 i 13.%
}

%%K t07-apar-1 apar
\VUsaut{4.0mm}
\VUcentre{13.2pt}{90}{SERIA DRUGA}
\VUsaut{3.0mm}
\VUfilet{16mm}
\VUsaut{5.6mm}
\VUtitre{15.0pt}{60}{TABLICA Nr 7}
\VUsaut{3.0mm}

%%K t07-tit-1 sub
\VUcentre{12.0pt}{0}{\textit{Scena pierwsza.}}
\VUsaut{3.0mm}

%%K t07-tit-2 sub
\VUpk{10.2pt}{40}{Wieś zimą.}
\VUsaut{4.0mm}

%%K t07-c1-01-1 p
1. --- Podczas ferii noworocznych towarzyszyłem ojcu do Nowego Miasta, wioski, gdzie
musiał załatwić kilka spraw handlowych.

%%K t07-c1-02-1 p
2. --- Jechaliśmy \VUgras{dyliżansem} \textsuperscript{(11)}, który kursuje między
miastem a tym miasteczkiem; ale ponieważ było bardzo zimno, ta podróż w mało wygodnym
pojeździe nie była zbyt przyjemna.

%%K t07-c1-03-1 p
3. --- Dlatego poczuliśmy prawdziwą przyjemność, kiedy zobaczyliśmy, jak
\VUgras{woźnica} zatrzymuje swój wóz na placu, przed \VUgras{zajazdem}
\textsuperscript{(2)} pod \textit{Białym Niedźwiedziem}. \VUgras{Oberżysta}
\textsuperscript{(4)} czekał na nas na progu swojego domu, spokojnie paląc swoją
\VUgras{fajkę}.

%%K t07-c1-03-2 p
Zaprosił nas do środka i nie dałem się prosić, żeby usiąść przed ogniem z chrustu, który
trzeszczał w palenisku. Wtedy bardzo drwiłem z ostrego \VUgras{wiatru północnego}, który
skrzypiał starym \VUgras{szyldem} \textsuperscript{(3)} i unosił czarnymi wirami
\VUgras{dym} \textsuperscript{(1)} wychodzący z komina.

%%K t07-c1-04-1 p
4. --- Była pora śniadania. Nie czekając dłużej, dobrze zjedliśmy prosty, ale smaczny
posiłek, który nam podano (*).

%%K t07-tit-3 sub
\VUpk{10.2pt}{40}{Kilka odwiedzin.}
\VUsaut{3.0mm}

%%K t07-c1-05-1 p
5. --- Po śniadaniu towarzyszyłem ojcu, który jest agentem ubezpieczeniowym, do jego
klientów. Najpierw odwiedziliśmy \VUgras{kowala} \textsuperscript{(5)}, który mieszka
obok zajazdu. Był zajęty podkuwaniem konia. Podziwiałem siłę jego pomocnika, który mocno
trzymał na swoim skórzanym fartuchu \VUgras{kopyto} konia, podczas gdy kowal umocowywał
\VUgras{podkowę} \textsuperscript{(6)} mocnymi uderzeniami młota.

%%K t07-c1-06-1 p
6. --- Potem poszliśmy do \VUgras{szewca} \textsuperscript{(7)} ze wsi, starego Jakuba.
Chociaż ma za szyld przepiękny \VUgras{but}, zadowala się przyszywaniem podeszew do
starych dziurawych butów.

%%K t07-c1-06-2 p
Starzec powiedział nam ze śmiechem: « W mieście byłem "szewcem" i nie miałem klientów.
Tutaj jestem "łataczem butów" i roboty jest dużo. »

%%K t07-c1-07-1 p
7. --- Mówił nam o swoim sąsiedzie \VUgras{golibrodzie} \textsuperscript{(8)}, z którego
klienci nie są zadowoleni. Jego \VUgras{brzytwy} są zawsze wyszczerbione, a jego
\VUgras{nożyce} nigdy dobrze nie strzygą.

Ale ponieważ jest jedynym golibrodą we wsi, trzeba rzecz jasna dać się golić i strzyc
jemu. Zresztą jest skąpym i nieuprzejmym człowiekiem.

%%K t07-c1-08-1 p
8. --- Na szczęście pozostali \VUgras{sąsiedzi} nie są do niego podobni. Na przykład
\VUgras{kołodziej} \textsuperscript{(9)} jest dobrym człowiekiem, pracowitym i
usłużnym. To przyjemność dać mu naprawić wóz. Jego robota jest zawsze wykończona.

%%K t07-c1-09-1 p
9. --- Stary Jakub nie skarżył się też na \VUgras{rymarza} \textsuperscript{(10)},
któremu czasem pomaga zszywać \VUgras{uprzęże}, kiedy robota jest pilna.

%%K t07-c1-09-2 p
Mój ojciec zapytał go o rodzinę: « Wszyscy mają się dobrze. Moja wnuczka jest w
\VUgras{szkole} \textsuperscript{(12)}. Jej \VUgras{nauczycielka} \textsuperscript{(13)}
jest z niej bardzo zadowolona, jest dobrą \VUgras{uczennicą} \textsuperscript{(14)}. Nie
jest podobna do mojego niedobrego syna; on myśli tylko o zabawie. \VUgras{Nauczyciel}
\textsuperscript{(15)} niczego nie zdoła go nauczyć. »

%%K t07-tit-4 sub
\VUsaut{3.0mm}
\VUpk{10.2pt}{40}{Wiejskie pogawędki.}
\VUsaut{3.0mm}

%%K t07-c1-10-1 p
10. --- Starzec opowiedział nam potem nowiny ze wsi. Ma się zbudować nową szkołę, bo ta
mieszcząca się w \VUgras{domu gminnym} stała się za mała. Zresztą to łatwe, bo wystarczy
kupić część ogrodu \VUgras{młynarza}. To będzie bardzo korzystne dla biedaka. Z powodu
zimna \VUgras{kamienie młyńskie} \VUgras{młyna} \textsuperscript{(16)} nie mogą już mleć
zboża, bo \VUgras{koło} \textsuperscript{(17)} zostało zatrzymane przez zamarznięty
\VUgras{lód} \textsuperscript{(25)} \VUgras{strumienia} \textsuperscript{(23)}.

%%K t07-c1-11-1 p
11. --- « Co do budowli, czy widział pan nasz nowy \VUgras{kościół}
\textsuperscript{(18)}? Jest bardzo malowniczy pod swoim śnieżnym płaszczem.
\VUgras{Wrony} \textsuperscript{(19)}, które nie poznają już swojej starej dzwonnicy,
smutno siadają na wielkim drzewie \VUgras{cmentarza} \textsuperscript{(20)}. »

%%K t07-c1-12-1 p
12. --- Ta myśl o cmentarzu przypomniała szewcowi osoby, które mój ojciec znał i które
umarły w zeszłym roku. « Ile \VUgras{grobów} \textsuperscript{(21)}, mój dobry panie,
przybyło do innych, pod \VUgras{cyprysem} \textsuperscript{(22)}! Śmierć skosiła naszego
starego notariusza. Wszyscy mieszkańcy wsi byli na jego \VUgras{pogrzebie}. »

%%K t07-c1-13-1 p
13. --- « Oto jego następca, który jeździ na łyżwach po strumieniu. Właśnie przyszedł
dać mi naprawić rzemień swojej \VUgras{łyżwy} \textsuperscript{(26)}, który się urwał. »

%%K t07-c1-14-1 p
14. --- « Oto także na brzegu strumienia nowy \VUgras{strażnik polny}
\textsuperscript{(27)}. To były żandarm, który przeraża urwisów ze wsi. Uciekają, gdy
tylko zobaczą jego \VUgras{getry} \textsuperscript{(29)} i jego \VUgras{kepi}
\textsuperscript{(28)}. »

%%K t07-c1-15-1 p
15. --- « Ha! oto stary Józef, który wiezie \VUgras{wózek z wiązkami drewna}
\textsuperscript{(30)} dla piekarza. --- To prawda, powiedział mój ojciec, poznaję jego
zaprzęg \VUgras{mułów} \textsuperscript{(31)}. Muszę z nim pomówić, skorzystam z
okazji. Do widzenia, ojcze Jakubie. »

%%K t07-c1-16-1 p
16. --- Właśnie kiedy wyszliśmy, dzieci ze wsi opuściły szkołę i biegnąc rzuciły się na
\VUgras{ślizgawkę} \textsuperscript{(33)} z ryzykiem upadku i złamania sobie ręki przy
\VUgras{upadku} \textsuperscript{(34)}. Jedno z nich wzięło swoje \VUgras{sanki}
\textsuperscript{(32)} u kołodzieja, posadziło w nich jednego z kolegów i ruszyło
biegiem.

%%K t07-c1-17-1 p
17. --- Inne rzuciły się na \VUgras{zaspę śnieżną} \textsuperscript{(37)} obok
\VUgras{mostu} \textsuperscript{(40)} i zaczęły bombardować \VUgras{bałwana}
\textsuperscript{(39)}, którego ulepiły poprzedniego dnia i któremu włożyły kapelusz,
fajkę i tornister przez ramię.

%%K t07-c1-18-1 p
18. --- Mało dbają o mroźny wiatr i zimno. Większość nie ma nawet \VUgras{szalika}
\textsuperscript{(35)}, a żeby się rozgrzać po walce na \VUgras{śnieżne kule}
\textsuperscript{(38)}, wystarczy im garść bardzo gorących \VUgras{kasztanów}
\textsuperscript{(42)} kupionych u starej \VUgras{przekupki} Jakubowej, która drży z
zimna pod swoją zbyt cienką \VUgras{chustą} \textsuperscript{(43)}.

%%K t07-tit-5 sub
\VUsaut{3.0mm}
\VUcentre{12.0pt}{0}{\textit{Scena druga.}}
\VUsaut{3.0mm}

%%K t07-tit-6 sub
\VUpk{10.2pt}{40}{Zagroda wiejska wiosną.}
\VUsaut{4.0mm}

%%K t07-c2-01-1 p
1. --- Mimo wszystkich przyjemności zimy z głęboką radością witamy powrót
\VUgras{wiosny}.

%%K t07-c2-01-2 p
Jakże radosny obraz przedstawia podwórze gospodarskie wieczorem, w miesiącu maju!

%%K t07-c2-02-1 p
2. --- Na widnokręgu \VUgras{słońce} \textsuperscript{(1)} powoli się kładzie, zalewając
\VUgras{wieś} \textsuperscript{(3)} swoimi purpurowymi \VUgras{promieniami}
\textsuperscript{(2)}. Pilny \VUgras{oracz} \textsuperscript{(6)} spieszy się zaorać
swoje \VUgras{pole} \textsuperscript{(4)}.

%%K t07-c2-02-2 p
Swoim \VUgras{pługiem} \textsuperscript{(5)} zaprzężonym do silnego \VUgras{konia}
\textsuperscript{(7)} kreśli \VUgras{bruzdę} \textsuperscript{(8)}, w którą
\VUgras{siewca} \textsuperscript{(9)} pełną garścią rzuci \VUgras{ziarno}, które wyda
złote kłosy.

%%K t07-c2-03-1 p
3. --- \VUgras{Kosiarze} \textsuperscript{(11)} zaopatrzeni w swoje kosy skosili trawę
na łące, przewracacze siana wysuszyli \VUgras{siano} \textsuperscript{(10)} obracając je
\VUgras{widłami} \textsuperscript{(12)} i grabiami, i oto wracają.

%%K t07-c2-03-2 p
Jedni idą przed ciężkim wozem ciągnionym przez woły; niosą swoje narzędzie na ramieniu.
Inni, bardziej leniwi, wygodnie ułożyli się na pachnącym sianie.

%%K t07-c2-04-1 p
4. --- Przechodząc pozdrawiają przyjaźnie \VUgras{ogrodnika} \textsuperscript{(16)}
zajętego w \VUgras{sadzie} \textsuperscript{(13)} przycinaniem \VUgras{drzew owocowych}
i szczepieniem \VUgras{grusz} \textsuperscript{(14)}. \VUgras{Śliwy}
\textsuperscript{(15)} zaczynają kwitnąć, a w dobrze nawiezionym \VUgras{ogrodzie
warzywnym} \textsuperscript{(17)} warzywa, kapusty i sałaty już się dobrze mają.

%%K t07-c2-04-2 p
Na murku piękny \VUgras{paw} \textsuperscript{(18)} rozkłada ogon, żeby dać podziwiać
jego przepiękne pióra.

%%K t07-tit-7 sub
\VUpk{10.2pt}{40}{Podwórze gospodarskie.}
\VUsaut{3.0mm}

%%K t07-c2-05-1 p
5. --- Drzwi \VUgras{stajni} \textsuperscript{(19)} są otwarte i widać drabinę na siano
i \VUgras{ściółkę} \textsuperscript{(23)} \VUgras{klaczy} \textsuperscript{(21)}, którą
\VUgras{stajenny} \textsuperscript{(20)} właśnie wyprzągł. \VUgras{Źrebię}
\textsuperscript{(22)}, tak zabawne ze swoimi długimi nogami, idzie za matką podskakując.

%%K t07-c2-06-1 p
6. --- Obok \VUgras{studni} \textsuperscript{(29)} \VUgras{krowy} \textsuperscript{(25)}
idą pić do drewnianego \VUgras{koryta} \textsuperscript{(26)}, które służy im za
poidło. \VUgras{Cielę} \textsuperscript{(24)} korzysta ze sposobności, żeby possać, a
\VUgras{byk} \textsuperscript{(27)}, wąchając dobry zapach paszy, radośnie ryczy i dumnie
podnosi głowę z potężnymi \VUgras{rogami} \textsuperscript{(28)}.

%%K t07-c2-07-1 p
7. --- Wszyscy pracują. \VUgras{Służąca} \textsuperscript{(32)} trzyma w ręku
\VUgras{sznur} \textsuperscript{(31)} studni. Czerpie wodę \VUgras{wiadrem}
\textsuperscript{(30)}, żeby napoić bydło. Obok \VUgras{stodoły} \textsuperscript{(33)}
mężczyzna grabi słomę, a przed drzwiami \VUgras{domu} \textsuperscript{(34)} stara
\VUgras{prządka} \textsuperscript{(35)} swoim kołowrotkiem i swoją \VUgras{kądzielą}
przędzie \VUgras{len} albo \VUgras{konopie} jak za dawnych czasów.

%%K t07-tit-8 sub
\VUsaut{3.0mm}
\VUpk{10.2pt}{40}{Gospodarz.}
\VUsaut{3.0mm}

%%K t07-c2-08-1 p
8. --- \VUgras{Gospodarz} \textsuperscript{(36)} kieruje wszystkimi tymi robotami i daje
wskazówki swoim parobkom. Zaleca schować pod dach \VUgras{bronę} \textsuperscript{(40)}
i \VUgras{motykę} \textsuperscript{(39)}, o których zapomniano obok \VUgras{budy}
\textsuperscript{(38)} \VUgras{psa} \textsuperscript{(37)}.

%%K t07-c2-09-1 p
9. --- Jego wołania płoszą \VUgras{gołębia} \textsuperscript{(41)}, który z całych
skrzydeł leci do \VUgras{gołębnika} \textsuperscript{(42)}, i \VUgras{kury}
\textsuperscript{(43)}, które spieszą wrócić do \VUgras{kurnika} \textsuperscript{(44)}.

%%K t07-c2-10-1 p
10. --- Przed \VUgras{owczarnią} \textsuperscript{(48)} oto stary \VUgras{pasterz}
\textsuperscript{(45)}, który przyprowadza z powrotem swoją \VUgras{trzodę}
\textsuperscript{(49)}. Trzymając swój \VUgras{kij pasterski} \textsuperscript{(46)},
którego nigdy mu nie brakuje, tak samo jak jego wypłowiałej \VUgras{bluzy}
\textsuperscript{(47)}, prowadzi do obory \VUgras{barany} \textsuperscript{(50)},
\VUgras{owce} \textsuperscript{(53)}, \VUgras{jagnięta} \textsuperscript{(54)},
\VUgras{kozy} \textsuperscript{(51)} i \VUgras{koźlęta} \textsuperscript{(52)}. Jego
wierny \VUgras{pies} \textsuperscript{(55)} Argus idzie ostatni i swoim szczekaniem
przynagla ociągających się.

%%K t07-c2-11-1 p
11. --- Cały ten hałas i ruch nie mącą spokoju dobrej \VUgras{mlecznej krowy}
\textsuperscript{(56)}, która pobrzękuje swoim \VUgras{dzwonkiem} \textsuperscript{(57)},
podczas gdy doi się ją przed drzwiami \VUgras{obory} \textsuperscript{(58)}.

%%K t07-c2-12-1 p
12. Zdaje się rozumieć, że musi dać swoje mleko, żeby \VUgras{gospodyni}
\textsuperscript{(60)} mogła przygotować \VUgras{masło} w \VUgras{maselnicy}
\textsuperscript{(61)} albo wyborne \VUgras{sery} \textsuperscript{(64)}, które ociekają
z deski położonej na \VUgras{kadzi} \textsuperscript{(63)}.

%%K t07-c2-13-1 p
13. --- Cała \VUgras{mleczarnia} \textsuperscript{(59)} jest utrzymywana bardzo czysto
przez \VUgras{parobka} \textsuperscript{(65)}, który właśnie umył \VUgras{garnki na
mleko} \textsuperscript{(62)} w wielkim wiadrze, które niesie w ręku.

%%K t07-tit-9 sub
\VUsaut{3.0mm}
\VUpk{10.2pt}{40}{Podwórze drobiu.}
\VUsaut{3.0mm}

%%K t07-c2-14-1 p
14. --- W swoich grubych drewniakach chodzi trochę ciężko i tak płoszy \VUgras{indyka}
\textsuperscript{(68)}, \VUgras{kaczki} \textsuperscript{(66)}, \VUgras{kaczory}
\textsuperscript{(67)} i \VUgras{gęsi} \textsuperscript{(69)}, które szeroko otwierają
\VUgras{swoje dzioby} \textsuperscript{(70)} i niespokojnie krzyczą.

%%K t07-c2-14-2 p
\VUgras{Kogut} \textsuperscript{(71)}, bardziej wojowniczy, podnosi swój czerwony
\VUgras{grzebień} \textsuperscript{(72)}, potrząsa piórami swojego \VUgras{ogona}
\textsuperscript{(73)} i wydaje dźwięczne « kukuryku ».

%%K t07-c2-15-1 p
15. --- \VUgras{Kwoka} \textsuperscript{(74)} grzebie na \VUgras{gnoju}
\textsuperscript{(81)} ze swoimi \VUgras{kurczętami} \textsuperscript{(75)}.
\VUgras{Prosię} \textsuperscript{(78)} ucieka do swojej matki, ogromnej
\VUgras{świni} \textsuperscript{(77)}, którą widzimy obok chlewu.

%%K t07-c2-16-1 p
16. --- W swoich przegrodach \VUgras{króliki} \textsuperscript{(79)} chciwie jedzą
delikatną \VUgras{trawę} \textsuperscript{(80)}, którą przynosi im córka gospodarza.
Janina bardzo lubi swoje króliki i codziennie, zebrawszy świeże \VUgras{jajka}
\textsuperscript{(76)} w kurniku, przychodzi odwiedzić swoich przyjaciół.
\VUsaut{6.0mm}
\VUfilet{22mm}
